\documentclass[conference]{IEEEtran}
\usepackage{amsmath,booktabs,array,url,graphicx,placeins}
\title{Data Preprocessing for Nigerian Sign Language Detection:\\Phase 2 Implementation and Strategy--Reality Audit}
\author{\IEEEauthorblockN{Candace Tariro Hunzwi}\\\IEEEauthorblockA{Carnegie Mellon University\\Engineering Artificial Intelligence}}
\begin{document}\maketitle
\begin{abstract}
This report audits an image-based preprocessing strategy for Nigerian Sign Language (NSL) sign detection using YOLO26. Prior exploratory analysis indexed 4,815 images from 137 class folders and 4,815 Pascal VOC XML files; 4,814 images yielded a geometrically valid extracted bounding box. Images were reported at $640\times480$ pixels, and class frequency ranged from four to 65 images (imbalance ratio 16.25). We translate the assignment's tabular requirements into technically valid controls for object detection, specify a reproducible image--annotation verification pipeline, and document the supervisor-directed baseline on the original, imbalanced images. Full-collection pipeline outputs and detector metrics are deliberately not presented as completed: the provided test output is a one-image smoke test, while the reported EDA values originate from the full initial audit. This separation prevents a successful code check from being misrepresented as a full-data experiment.
\end{abstract}
\begin{IEEEkeywords}Nigerian Sign Language, object detection, preprocessing, YOLO26, annotation audit, leakage, reproducibility\end{IEEEkeywords}
\section{Problem and Data}
The NSL dataset is a multi-class object-detection dataset, not a tabular predictor matrix. Each observation is an RGB image with one or more associated sign-class bounding boxes. The original study used an earlier YOLOv5 model~\cite{kolawole}; the current project targets YOLO26. An initial audit of the inspected image-folder/Pascal VOC collection reported the values in Table~\ref{tab:audit}. These are prior notebook audit outputs and are distinct from an independent full re-execution of the accompanying Phase 2 program.
\begin{table}[!b]\caption{Previously Reported Raw-Data Audit}\label{tab:audit}\centering\begin{tabular}{lr}\toprule Quantity & Value\\\midrule Indexed images & 4,815\\XML files & 4,815\\Folder-derived sign classes & 137\\Images with valid extracted box & 4,814\\Raw image dimensions & $640\times480$\\Images per class (min--max) & 4--65\\Largest/smallest class ratio & 16.25\\\bottomrule\end{tabular}\end{table}
\FloatBarrier
The discrepancy between XML-file count and usable-box count must be resolved by inspecting the affected file; a matched XML is not proof of valid supervision. The original class list and XML label names also require consistency checks before training.
\section{Diagnostic Evidence and Decisions}
\subsection{Target Imbalance}
The class distribution has a ratio $65/4=16.25$. This is evidence of uneven representation, not evidence by itself of detector bias. In accordance with supervisor guidance, the baseline retains the natural class frequencies; targeted balancing is deferred until the baseline is evaluated. Training-only augmentation or sampling can then be tested without altering held-out data. The class-distribution figure is available in the project README under Fig. 1 in the accompanying repository.
\subsection{Missingness and Annotation Integrity}
For object detection, missing or invalid supervision is a data-integrity problem rather than a value to impute. The pipeline verifies that each image is readable, has exactly one parsable XML annotation, uses a known class label, and contains finite, positive bounding-box coordinates whose edges lie within the image. It also flags empty annotations and image--annotation mismatches. A missing image, sign label, or bounding box is never replaced with a mean, median, or predicted value; unresolved records are excluded from every split and recorded in an exclusion log with the file identifier and reason. These checks establish structural validity only. They do not confirm that a geometrically valid box encloses the intended sign, so flagged and sampled accepted annotations still require visual review.
\subsection{Outliers and Image Geometry}
Rare or extreme bounding-box areas can represent valid examples; box-size outliers should therefore be flagged and reviewed rather than capped or deleted automatically. The raw $640\times480$ images are preserved in the baseline, and the training loader letterboxes to a target $640\times640$ canvas while retaining aspect ratio. At that exact input scale, 80 pixels of vertical padding may be used on each side, although optimized rectangular batching can produce other intermediate shapes. Ground-truth boxes must remain aligned with image geometry.
\subsection{Background and Visual Confounding}
Representative images show additional people and classroom backgrounds. These observations motivate, but do not validate, future signer-background-removal and padded sign-region-cropping studies. Cropping a test image from its \emph{ground-truth} box would use information unavailable in deployment; such results must be labeled oracle-assisted and not compared as ordinary end-to-end detections.
\section{Leakage-Free Implementation}
The accompanying diagnostics notebook contains raw-data indexing and visualization code. A separate, reusable Python preprocessing module reads source images and XML annotations, validates and converts bounding boxes into normalized YOLO coordinates, creates a deterministic split manifest, and records exact-image SHA-256 identifiers. Exact duplicate hashes are assigned to the same split. This does not establish signer-disjointness or catch all near duplicates; both are outstanding audit requirements. Class frequencies are not equalized. Random transformations, if enabled in future training, must occur only on training examples.
\par For rubric alignment, a Scikit-Learn \texttt{Pipeline} containing a stateless \texttt{ImageAuditFeatures} transformer is provided for reproducible audit metadata (width, height, aspect ratio, intensity mean and standard deviation). Its \texttt{fit} is invoked on training paths only. This yields a five-column \emph{audit} matrix, not YOLO's model input. A \texttt{ColumnTransformer} for pixel values is not introduced because image tensor preprocessing and box-aware augmentation belong in the detection pipeline; there is no tabular predictor block requiring fitted numeric/categorical transforms. This is a scope adaptation, not a claim that YOLO alone fulfills the assignment's explicit \texttt{ColumnTransformer} requirement; instructor confirmation is needed.

\subsection{Run Contract and Evidence}
Before a detector is trained, the preprocessing program must be run against the complete local dataset with a fixed seed. It must write a validated manifest, exclusion log, split manifest, class-map YAML, verification JSON, and 300-DPI diagnostic figures. Each excluded record requires an image identifier, reason, and action. The validation and test partitions must never determine conversion rules, augmentation policy, hyperparameters, or the confidence threshold. The current supplied \texttt{test\_outputs} directory is useful only as a smoke test: it indexes one image, produces a $1\times5$ audit matrix, and has empty validation/test partitions. It is not evidence for the full collection.
\section{Strategy vs. Reality Decision Log}
Table~\ref{tab:decision} records the relationship between the Phase~1 plan, the evidence observed so far, and the work that remains. It separates implemented checks from proposed experiments so that planned work is not reported as a result.
\begin{table}[!t]
\caption{Strategy--Reality Decision Log}
\label{tab:decision}
\centering
\scriptsize
\begin{tabular}{p{0.22\columnwidth}p{0.68\columnwidth}}
\toprule
Area & Evidence and decision \\
\midrule
Target balance & 137 classes with 4--65 images (16.25:1 ratio). The baseline retains the original distribution; balancing is deferred. \\
Annotation failures & 4,815 XML files and 4,814 valid extracted boxes were reported. Validation/exclusion code exists; full-data rerun and manual review remain pending. \\
Geometry & Raw images are $640\times480$; the target canvas is $640\times640$. YOLO-coordinate conversion is implemented; runtime loader verification is pending. \\
Outliers & Area distributions are reported in the README. Extreme but valid cases are retained for review rather than removed automatically. \\
Leakage & Signer identity is unknown and duplicates are possible. Exact-hash grouping is scripted; perceptual-duplicate and signer checks remain pending. \\
Architecture & YOLO26 is selected as the baseline; training outcomes are pending. \\
Background & The original-image baseline is retained. Background removal and cropping are future experiments. \\
\bottomrule
\end{tabular}
\end{table}
\FloatBarrier
\section{YOLO26 Selection, Hyperparameters, and Baseline}
YOLO26 supports bounding-box detection and pretrained checkpoint fine-tuning, providing a practical route from the original YOLOv5-based sign-to-speech design. Published COCO results for YOLO26n report 2.4 million fused-model parameters, 40.9 mAP$_{50:95}$ and 1.7 ms inference using T4 TensorRT under the documented benchmark configuration~\cite{yolo26}. YOLO26s reports 9.5 million parameters, 48.6 mAP$_{50:95}$ and 2.5 ms on the same benchmark. These are not performance claims for NSL. The Nano variant is a defensible initial efficiency-oriented baseline; Small may later be considered if validation results justify additional capacity.
\par The proposed training configuration uses pretrained \texttt{yolo26n.pt}, $640$-pixel input, initial budget of 100 epochs, nominal batch size 16, seed 42, and a fixed 70/15/15 target split. Actual split sizes, model arguments, accelerator, package versions, any default augmentation, training logs and best validation metrics must be recorded after execution. The final test set should be used only once model selection is frozen.
\section{Output Verification and Limitations}
Table~\ref{tab:verify} separates prior full-collection EDA from the values that the final preprocessing execution must measure. The full-dataset program writes \texttt{verification.json} with indexed, accepted and excluded counts, exact train/validation/test image counts, a Boolean exact-hash leakage check, and shape $(N_{\mathrm{train}},5)$ for the auxiliary audit feature matrix. The detector input has expected shape $(N,3,640,640)$ after letterboxing when full square batches are used; this requires runtime inspection and is not claimed as measured. The script also exports 300-DPI validated class-count and box-area plots. Full-collection values cannot be inserted until the program is run against the 4,815 original images.
\begin{table}[t]
\caption{Evidence Status for Pipeline Verification}\label{tab:verify}
\centering\footnotesize
\begin{tabular}{p{0.47\columnwidth}p{0.42\columnwidth}}
\toprule
Evidence item & Status in this report\\\midrule
Raw image/XML inventory & Prior full-collection EDA: 4,815 / 4,815\\
Valid extracted boxes & Prior full-collection EDA: 4,814; one anomaly remains\\
Audit matrix dimensions & Smoke test only: $1\times5$; full run pending\\
Train/validation/test counts & Pending full-data split\\
Exact-hash leakage check & Smoke test passes; full-data result pending\\
Signer-disjoint verification & Not established; signer metadata unavailable\\
YOLO tensor shape & Expected after letterbox; runtime confirmation pending\\
Detector metrics & Not run; no NSL metric claimed\\
\bottomrule
\end{tabular}
\end{table}
\subsection{Notebook--Pipeline Reconciliation}
The uploaded diagnostic notebook contains a separate baseline-preparation draft; the revised notebook removes duplicated path-dependent conversion cells and directs readers to the modular Python scripts. The baseline must be trained only after the generated class map is checked against the authoritative 137-label YOLO map. Stratified image-level splitting is insufficient to establish unseen-signer generalization when signer identifiers are absent. Exact-hash grouping is a partial leakage control, not a proof of full independence.
\par The repository currently contains raw-data class-distribution, annotation-coverage, image-dimension, box-distribution, representative-image, and transformation figures. These support the Phase~1-to-Phase~2 rationale but are not detector-performance plots. The final report should embed or explicitly reference only figures generated from the final pipeline run, at 300 DPI. Missingness heatmaps, raw-pixel correlation/VIF plots, and PCA scree plots require an agreed image-appropriate analogue; they must not be fabricated merely to imitate a tabular rubric.
\section{Conclusion}
The confirmed EDA motivates geometry-preserving preprocessing, annotation-integrity checks, duplicate-aware splitting, and explicit uncertainty about rare classes. The immediate controlled experiment is YOLO26 on the original class distribution and full images, not a balanced or background-segmented derivative. The notebook preserves implemented data-indexing and diagnostic cells; the full-data preprocessing script and baseline training are separate stages. The present artifacts distinguish prior EDA, a one-image code smoke test, and still-unexecuted full-data/model stages. No gains are claimed before evaluation.
\begin{thebibliography}{2}
\bibitem{kolawole} S. Kolawole et al., ``Sign-to-Speech Model for Sign Language Understanding: A Case Study of Nigerian Sign Language,'' 2021. [Online]. Available: \url{https://arxiv.org/abs/2111.00995}
\bibitem{yolo26} Ultralytics, ``YOLO26 Documentation,'' 2026. [Online]. Available: \url{https://docs.ultralytics.com/models/yolo26/}
\end{thebibliography}
\end{document}
